\documentclass[11pt]{article}
\usepackage[a4paper,margin=26mm]{geometry}
\usepackage{amsmath,amssymb,mathtools}
\usepackage{unicode-math}
\setmainfont{Latin Modern Roman}
\setmathfont{Latin Modern Math}
\usepackage{xurl,hyperref}
\hypersetup{hidelinks,pdftitle={UnsolvedMath Problem 10300038}}
\newcommand{\sref}[2]{\hyperlink{source:#1:#2}{[S#2]}}
\newcommand{\cataloguescope}[1]{#1}
\begin{document}
% UnsolvedMath ID 10300038; source catalogue code AMR-102-0038
\setcounter{section}{10300037}
\section{Classical 3-manifold theory — Question 9.5}\label{10300038}
{\small\sffamily UnsolvedMath ID 10300038 · Topology}
\subsection{Definitions and mathematical statement}

\paragraph{Shared notation.}$\mathbb N=\{1,2,\ldots\}$, and $\mathbb Z,\mathbb Q,\mathbb R,\mathbb C$ denote the integers, rationals, reals and complex numbers. Any other conventions are specified locally.


A surface lamination $\Lambda$ is a closed subset of $M$ locally homeomorphic to a product of a disk with a closed transverse set. It is essential when it has no sphere leaf or torus leaf bounding a solid torus, and its complementary regions are irreducible, have incompressible boundary, and have no compressing monogon. Such a monogon is a compressing disk with one boundary arc on a boundary leaf and the other along a transverse interval fiber. It is genuine if some complementary region is not an interval bundle. A sublamination is a closed union of leaves. A Heegaard splitting $M=H_1\cup_\Sigma H_2$ has distance equal to the distance in the curve graph of $\Sigma$ between the meridian-disk boundary sets of $H_1,H_2$. The curve graph has essential simple closed curves as vertices and disjointness edges.
\paragraph{Statement recorded in the supplied source.}
Determine what useful constraints or filtration structure Heegaard-splitting distance imposes on essential laminations in $M$. In particular, if $M$ has a Heegaard splitting of distance at least $2g$, what conclusions can be drawn about its essential laminations? Compare the surface theorem that an essential surface of genus $g$ bounds splitting distance above by $2g$, while seeking genuinely lamination-theoretic consequences. \sref{10300038}{2}
\subsection{Short English statement}
Determine how large Heegaard distance constrains essential laminations, extending the familiar distance bound from incompressible surfaces.
\subsection{Sources}
\begin{description}
\item[\hypertarget{source:10300038:1}{S1}] ULAM.AI and the UnsolvedMath Contributors, \emph{UnsolvedMath: A Curated Collection of Open Mathematics Problems}, record 10300038 (AMR-102-0038). CC BY 4.0. \url{https://huggingface.co/datasets/ulamai/UnsolvedMath}.
\item[\hypertarget{source:10300038:2}{S2}] Danny Calegari, \emph{Problems in foliations and laminations of 3-manifolds}, version 0.78 (2002), Question 9.5; Definitions 1.1--1.10 and local definitions and remarks. \url{https://arxiv.org/abs/math/0209081v1}.
\end{description}
\label{end:10300038}
\end{document}
